%============================================================
\chapter{Đối Nhân Xử Thế (The Art of Dealing with People)}
\begin{center}
  \textit{\color{truyenblue}You can make more friends in two months by becoming\\
  interested in other people than in two years by\\
  trying to get other people interested in you.}\\
  \textit{(Bạn có thể kết bạn trong hai tháng bằng cách quan tâm đến người khác\\
  hơn là trong hai năm cố làm người khác quan tâm đến bạn.)}\\[4pt]
  \small\color{truyengold}--- Dale Carnegie ---
\end{center}
\ngancach

\section{Gia Cát Lượng Và Bài Học Dùng Người}
\begin{truyen}{Gia Cát Lượng -- Dùng Người Đúng Chỗ}{Tam Quốc Chí, Trung Hoa}
\chuhoa{G}{}G ia Cát Lượng -- mưu sĩ thiên tài của Thục Hán -- nổi tiếng với nghệ thuật dùng người.\\
\textit{(Zhuge Liang, the brilliant strategist of Shu Han, was famous for the art of utilizing people.)}

Mã Siêu dũng mãnh nhưng nóng nảy -- ông dùng ở tiên phong xung trận, không dùng lập kế.\\
\textit{(Ma Chao was brave but hot-headed -- he used him at the vanguard, never for planning.)}

Triệu Vân trung thành cẩn thận -- ông giao việc bảo vệ hậu phương và thân vương.\\
\textit{(Zhao Yun was loyal and careful -- he assigned him to protect the rear and the princes.)}

Khương Duy thông minh nhạy bén -- ông truyền nghề và chọn làm người kế thừa chiến lược.\\
\textit{(Jiang Wei was clever and perceptive -- he mentored him and chose him as his strategic successor.)}

Ngay cả Mã Tắc -- người ông yêu quý nhất -- khi dùng sai việc ông vẫn xử theo quân pháp.\\
\textit{(Even Ma Su -- the one he loved most -- when misused in the wrong task, he still executed him according to military law.)}

Ông nói: ``Dùng người như dùng khí: phát huy sở trường, tránh sở đoản, kết quả mới tốt được.''\\
\textit{(He said: ``Using people is like using tools: leverage strengths, avoid weaknesses, only then will the results be good.'')}

\end{truyen}
\begin{baihoc}
  Nghệ thuật đối nhân xử thế cao nhất là nhận ra điểm mạnh trong mỗi con người và giúp họ tỏa sáng đúng chỗ.\\
  \textit{(The highest art of dealing with people is recognizing the strengths in each person and helping them shine in the right place.)}
\end{baihoc}

\section{Abraham Lincoln Và Kẻ Thù Trở Thành Bạn}
\begin{truyen}{Lincoln -- Phá Hủy Kẻ Thù Bằng Cách Biến Họ Thành Bạn}{Nước Mỹ, Thế Kỷ XIX}
\chuhoa{M}{}M ột người bạn hỏi Lincoln: ``Tại sao ông không tiêu diệt kẻ thù của ông đi?''\\
\textit{(A friend asked Lincoln: ``Why don't you destroy your enemies?'')}

Lincoln mỉm cười: ``Tôi đã tiêu diệt họ rồi -- bằng cách biến họ thành bạn bè của tôi.''\\
\textit{(Lincoln smiled: ``I have already destroyed them -- by making them my friends.'')}

Edwin Stanton từng công khai lăng mạ Lincoln, gọi ông là ``con khỉ đen bẩn thỉu.''\\
\textit{(Edwin Stanton had publicly insulted Lincoln, calling him a ``filthy black ape.'')}

Lincoln không trả thù -- ông bổ nhiệm Stanton làm Bộ trưởng Quốc phòng vì tài năng của ông ta.\\
\textit{(Lincoln did not retaliate -- he appointed Stanton as Secretary of War because of the man's talent.)}

Stanton trở thành người hết lòng phục vụ Lincoln và sau khi Lincoln bị ám sát, ông đứng bên quan tài khóc: ``Đây là người vĩ đại nhất tôi từng biết.''\\
\textit{(Stanton became Lincoln's most devoted servant and when Lincoln was assassinated, he stood by the coffin weeping: ``He was the greatest man I have ever known.'')}

\end{truyen}
\begin{baihoc}
  Kẻ thù giỏi nhất không phải là người đánh bại đối thủ mà là người chuyển hóa đối thủ thành đồng minh.\\
  \textit{(The best kind of enemy-handling is not defeating opponents but transforming them into allies.)}
\end{baihoc}

\section{Bác Hồ Và Nghệ Thuật Lắng Nghe}
\begin{truyen}{Bác Hồ -- Lãnh Tụ Lắng Nghe}{Việt Nam, Thế Kỷ XX}
\chuhoa{Đ}{}Đ iều đặc biệt của Hồ Chí Minh là khả năng lắng nghe -- từ cụ nông dân đến em thiếu nhi.\\
\textit{(What was special about Ho Chi Minh was his capacity to listen -- from elderly farmers to young children.)}

Khi thăm làng quê, ông ngồi xuống đất cùng dân, hỏi thăm từng gia đình, không ngại mùi bùn đất.\\
\textit{(When visiting villages, he sat on the ground with the people, inquired about each family, unmindful of the smell of mud and earth.)}

Mỗi buổi tiếp dân, ông ghi chép tỉ mỉ và theo dõi đến khi vấn đề được giải quyết.\\
\textit{(At each audience with the people, he took meticulous notes and followed up until each issue was resolved.)}

Học trò hỏi bí quyết thu phục lòng người, ông đáp: ``Không có bí quyết. Hãy thực sự quan tâm đến họ.''\\
\textit{(When students asked the secret to winning people's hearts, he replied: ``There is no secret. Truly care about them.'')}

\end{truyen}
\begin{baihoc}
  Người có tầm ảnh hưởng lớn nhất không phải người nói nhiều nhất mà là người lắng nghe sâu nhất.\\
  \textit{(The most influential person is not the one who talks the most but the one who listens the deepest.)}
\end{baihoc}

\section{Khổng Minh Thả Mạnh Hoạch Bảy Lần}
\begin{truyen}{Thất Cầm Mạnh Hoạch -- Nghệ Thuật Chinh Phục Lòng Người}{Tam Quốc}
\chuhoa{K}{}K hi đánh dẹp phương Nam, Gia Cát Lượng bắt được thủ lĩnh Mạnh Hoạch nhưng không giết.\\
\textit{(When pacifying the south, Zhuge Liang captured chieftain Meng Huo but did not execute him.)}

Ông thả Mạnh Hoạch về và cho binh lính tập hợp lại. Mạnh Hoạch tức tối kéo quân quay lại đánh.\\
\textit{(He released Meng Huo and let him regroup his troops. Meng Huo furiously returned to attack.)}

Cứ như vậy, bắt rồi thả, thả rồi bắt -- bảy lần liên tiếp.\\
\textit{(This went on, captured then released, released then recaptured -- seven consecutive times.)}

Đến lần thứ bảy, Mạnh Hoạch quỳ xuống: ``Thừa tướng tha bảy lần bắt bảy lần -- tôi biết thần oai của ngài, xin quy phục.''\\
\textit{(On the seventh time, Meng Huo knelt: ``The Prime Minister has released and recaptured me seven times -- I know your divine authority, I submit.'')}

Gia Cát Lượng chinh phục phương Nam không bằng gươm giáo mà bằng tâm phục khẩu phục.\\
\textit{(Zhuge Liang pacified the south not by the sword but by winning genuine submission of heart and mouth.)}

\end{truyen}
\begin{baihoc}
  Chinh phục bằng sức mạnh tạo ra kẻ phục tùng; chinh phục bằng nhân tâm tạo ra đồng minh chân thành.\\
  \textit{(Conquering by force creates subjects; conquering by humanity creates genuine allies.)}
\end{baihoc}

\section{Dale Carnegie Và Cái Tên}
\begin{truyen}{Dale Carnegie -- Phép Màu Của Cái Tên}{Nước Mỹ, Thế Kỷ XX}
\chuhoa{D}{}D ale Carnegie -- tác giả ``Đắc Nhân Tâm'' -- nhận thấy điều đơn giản nhất mà người ta thường bỏ qua.\\
\textit{(Dale Carnegie, author of ``How to Win Friends and Influence People,'' noticed the simplest thing people often overlook.)}

Ông viết: ``Gọi đúng tên của một người -- đó là âm thanh ngọt ngào nhất đối với tai họ.''\\
\textit{(He wrote: ``A person's name -- that is the sweetest sound in any language to that person's ears.'')}

Ông kể về tổng thống Roosevelt: mỗi người ông gặp dù là người hầu hay quan chức, ông đều nhớ tên.\\
\textit{(He told of President Roosevelt: every person he met, whether a servant or an official, he remembered by name.)}

Roosevelt giải thích: ``Nhớ tên người khác là cách đơn giản nhất để tỏ ra rằng bạn quý trọng họ.''\\
\textit{(Roosevelt explained: ``Remembering someone's name is the simplest way to show you value them.'')}

Carnegie kết luận: ``Những kỹ năng phức tạp nhất trong đối nhân xử thế thường đến từ những điều đơn giản nhất.''\\
\textit{(Carnegie concluded: ``The most complex interpersonal skills often come from the simplest things.'')}

\end{truyen}
\begin{baihoc}
  Quan tâm thực sự đến người khác -- bắt đầu từ việc nhớ tên họ -- là nền tảng của mọi mối quan hệ tốt.\\
  \textit{(Genuinely caring about others -- starting with remembering their name -- is the foundation of every good relationship.)}
\end{baihoc}
